% 10-roadmap.tex
\chapter{Future Work and Roadmap}
\label{ch:roadmap}

This chapter outlines the planned development trajectory for Doki
across three time horizons. Each milestone is driven by user feedback,
emerging platform capabilities, and the evolving container ecosystem.

\section{Short-term (v0.10)}
\label{sec:roadmap-short}

The v0.10 release focuses on closing functional gaps and improving
the developer experience for existing features.

\begin{itemize}[leftmargin=*, itemsep=4pt]
  \item \textbf{Health check execution.} Implement runtime execution
        of health check commands for Compose \texttt{service\_healthy}
        conditions. Currently, the condition is parsed but not executed
        by the container runtime. This will enable proper readiness and
        liveness probes in multi-service deployments.

  \item \textbf{Certificate rotation tooling.} Automate the rotation
        of DokiLink CA and link certificates. CA lifetime is currently
        365 days; link certificates expire after 90 days. A built-in
        rotation command will prevent certificate expiration from
        disrupting mesh connectivity.

  \item \textbf{Overlay2 on rootless Android.} Investigate kernel
        overlay support on rooted Android devices to improve storage
        performance beyond fuse-overlayfs. Native overlay2 would reduce
        I/O overhead by approximately 10\% based on preliminary
        benchmarks.

  \item \textbf{Compose profiles UI.} Improve the Compose profile
        experience with explicit profile activation feedback,
        dependency visualization, and a \texttt{doki compose profiles}
        command that lists available profiles and their services.
\end{itemize}

\section{Medium-term (v0.11)}
\label{sec:roadmap-medium}

The v0.11 release introduces advanced security features and protocol
upgrades.

\begin{itemize}[leftmargin=*, itemsep=4pt]
  \item \textbf{Noise protocol (L3).} Implement the Noise protocol
        framework~\cite{noise2015} for DokiLink encryption,
        providing forward secrecy and post-quantum resistance. This
        replaces the current static key exchange with ephemeral
        Diffie-Hellman.

  \item \textbf{Wire protocol upgrade.} Replace the JSON gossip
        protocol with gRPC/protobuf for improved performance and
        type safety. Benchmarking shows protobuf serialization is
        3--5x faster than JSON for the peer discovery message format.

  \item \textbf{gVisor integration.} Enable the gVisor user-space
        kernel~\cite{gvisor2024} as an isolation option on
        Android devices with sufficient resources. gVisor provides
        stronger syscall filtering than proot while remaining
        rootless.

  \item \textbf{Image signing.} Implement OCI image signature
        verification using Sigstore/cosign integration. This enables
        supply chain security for images pulled from public registries.
\end{itemize}

% Shipped ahead of schedule in v0.11.0 — see notes below.
\section{Long-term (v1.0)}
\label{sec:roadmap-long}

The v1.0 release targets production-grade features and ecosystem
integration.

\begin{itemize}[leftmargin=*, itemsep=4pt]
  \item \textbf{WireGuard mesh.} Full WireGuard integration for
        encrypted point-to-point tunnels between Doki instances.
        WireGuard provides kernel-level encryption with minimal
        overhead compared to userspace TLS.
\end{itemize}

% Implemented ahead of schedule in v0.11.0.
\subsection*{Shipped in v0.11.0 (ahead of schedule)}

The following features were originally planned for v1.0 but were
implemented ahead of schedule in the v0.11.0 release:

\begin{itemize}[leftmargin=*, itemsep=4pt]
  \item \textbf{NAT traversal.} STUN/TURN hole-punching for peer
        discovery across NAT boundaries, enabling DokiLink
        connectivity between devices on different networks without
        static IP addresses.

  \item \textbf{DHT peer discovery.} Distributed hash table for
        decentralized peer discovery without static configuration,
        eliminating the need for manually adding peers via
        \texttt{doki link add}.

  \item \textbf{Kubernetes compatibility.} Full CRI (Container
        Runtime Interface) compliance for Kubernetes node integration,
        enabling Doki to serve as a runtime for Kubernetes pods on
        Android devices.
\end{itemize}

\section{Feature Priority Matrix}
\label{sec:priority-matrix}

The following table summarizes planned features with their estimated
impact and implementation complexity.

\begin{table}[htbp]
  \centering
  \caption{Feature priority matrix for the Doki roadmap.}
  \label{tab:roadmap-priorities}
  \begin{tabularx}{\textwidth}{l l c c X}
    \toprule
    \textbf{Feature} & \textbf{Version} & \textbf{Impact} & \textbf{Complexity} & \textbf{Rationale} \\
    \midrule
    Health checks     & v0.10 & High   & Low    & Unblocks Compose service\_healthy \\
    Cert rotation     & v0.10 & Medium & Low    & Prevents mesh connectivity loss \\
    Overlay2          & v0.10 & Medium & High   & Requires kernel investigation \\
    Noise protocol    & v0.11 & High   & Medium & Forward secrecy for DokiLink \\
    gVisor            & v0.11 & High   & High   & Strongest rootless isolation \\
    Image signing     & v0.11 & Medium & Medium & Supply chain security \\
    WireGuard mesh    & v1.0  & High   & High   & Production-grade networking \\
    NAT traversal     & v1.0  & High   & High   & Cross-network connectivity \\
    K8s CRI           & v1.0  & High   & Very High & Ecosystem integration \\
    \bottomrule
  \end{tabularx}
\end{table}

\section{Development Timeline}
\label{sec:timeline}

\begin{figure}[htbp]
  \centering
  \begin{tikzpicture}[
    node distance=0.5cm,
    every node/.style={font=\sffamily\footnotesize},
    milestone/.style={
      rectangle, rounded corners=3pt,
      draw=nodeblue!60, fill=fillblue,
      minimum width=2.2cm, minimum height=0.7cm,
      text=textdark, align=center, inner sep=4pt
    },
    phase/.style={
      rectangle, rounded corners=2pt,
      draw=nodeorange!50, fill=fillorange!30,
      minimum width=1.8cm, minimum height=0.6cm,
      text=textdark, align=center, inner sep=3pt,
      font=\sffamily\scriptsize\bfseries
    }
  ]
    % Timeline axis
    \draw[linegray, thick] (0,0) -- (12,0);

    % Phase markers
    \node[phase] at (1.5, 0.8) {v0.9.3};
    \node[phase] at (4.5, 0.8) {v0.10};
    \node[phase] at (8, 0.8) {v0.11};
    \node[phase] at (11, 0.8) {v1.0};

    % Current
    \fill[nodeblue] (1.5,0) circle (3pt);
    \node[below=0.3cm, font=\sffamily\tiny\bfseries, text=nodeblue] at (1.5,0) {Current};

    % Features
    \node[milestone, below=0.6cm] at (4.5,0) (h1) {Health checks};
    \node[milestone, below=0.6cm] at (8,0) (n1) {Noise protocol};
    \node[milestone, below=0.6cm] at (11,0) (w1) {WireGuard mesh};

    \draw[->, nodeorange, dashed, line width=0.6pt] (1.5,0) -- (4.5,0);
    \draw[->, nodeorange, dashed, line width=0.6pt] (4.5,0) -- (8,0);
    \draw[->, nodeorange, dashed, line width=0.6pt] (8,0) -- (11,0);
  \end{tikzpicture}
  \caption{Doki development roadmap timeline with key milestones.}
  \label{fig:roadmap-timeline}
\end{figure}
